\documentclass[11pt]{article}

\usepackage[T1]{fontenc}
\usepackage[margin=1in]{geometry}
\usepackage{amsmath,amssymb,mathtools,amsthm}
\usepackage{booktabs,tabularx,array}
\IfFileExists{lmodern.sty}
  {\usepackage{lmodern}\usepackage{microtype}}
  {\usepackage[expansion=false]{microtype}}
\usepackage{enumitem}
\usepackage[hidelinks]{hyperref}

\setlength{\parindent}{0pt}
\setlength{\parskip}{0.55em}
\setlist[itemize]{leftmargin=*,itemsep=0.15em,topsep=0.25em}
\renewcommand{\arraystretch}{1.12}

\newcommand{\AP}{\operatorname{AP}}
\newcommand{\CNAP}{\operatorname{CNAP}}
\newcommand{\Regime}{\mathcal{R}}
\newcommand{\Assess}{\mathcal{A}}
\newcommand{\CSet}{\mathcal{C}}
\newcommand{\PSet}{\Pi_\star}
\newcommand{\E}{\mathbb{E}}
\newcommand{\Prb}{\mathbb{P}}
\newcommand{\iid}{\mathrel{\overset{\mathrm{iid}}{\sim}}}
\newcommand{\APta}{\widetilde{\AP}}
\newcommand{\CNAPta}{\widetilde{\CNAP}}
\newcommand{\Dta}{\widetilde{\Delta}}
\newcommand{\RobDta}{\underline{\widetilde{\Delta}}_{\PSet}}

\newtheorem{definition}{Definition}
\newtheorem{proposition}{Proposition}
\theoremstyle{remark}
\newtheorem*{remark}{Remark}

\title{Supported Model Replacement Across Target Prevalences\\
\large A Paired Replication-Based Method for Average Precision}
\author{Marcus A. Thomas}
\date{Methodological proposal}

\begin{document}

\maketitle

\begin{abstract}
Average precision depends on the prevalence of the population being evaluated, while a model-replacement claim may fail when the evaluation is repeated. A defensible comparison must therefore preserve pairing while challenging the claimed advantage across both target populations and independent replications. Let $\PSet$ be a declared set of target prevalences, let $\Regime$ be a law for one evaluation containing the outcomes and both models' scores, and let $K$ be the order of the replication challenge. For candidate model $A$ and incumbent model $B$, prior-standardize average precision to prevalence $\pi$ and normalize population random ranking to zero. The proposed replacement estimand is
\[
\theta_{A:B}
=
\E_{\Regime}\left[
\min_{1\leq j\leq K}
\inf_{\pi\in\PSet}
\left\{
\widetilde{\CNAP}_{A,\pi}(D_j)
-
\widetilde{\CNAP}_{B,\pi}(D_j)
\right\}
\right].
\]
The difference is formed on the same units before either challenge is applied, so favorable conditions for one model cannot be matched against unfavorable conditions for the other. The tildes denote average precision averaged over the orderings within each score-tie block. This convention removes expected credit for arbitrary score refinement and, under an exchangeable-label reference law, prevents differing tie structures from creating a positive supported advantage in either orientation. With several empirical replications, a complete U-statistic estimates the criterion directly. With one evaluation, a paired stratified bootstrap provides only projected support under an empirical fixed-model regime. The estimand is an integrated joint-survival magnitude, not the probability of clearing a fixed margin. Replacement requires a validated lower uncertainty bound for $\theta_{A:B}$ to exceed a prespecified practical margin; literal survival of that margin is reported separately.
\end{abstract}


\section{Performance Claims and their Support}
\label{sec:introduction}

Prediction systems often rank cases so that a selected set can receive further
attention. A radiology service, for example, may send its highest-scored images
for review. The composition of that set matters, not only the separation
between outcome classes, because the selected set is what the decision maker
receives.

Separation (eg AUROC) and selected-set purity (eg Average Precision) part company when prevalence changes.
Suppose a threshold selects $80\%$ of positive cases and $5\%$ of negative
cases. In a balanced population of 10,000 cases, it returns 4,000 positives and
250 negatives, for precision $0.94$. In a population with $1\%$ positives, the
same threshold returns 80 positives and 495 negatives, for precision $0.14$.
Its true-positive and false-positive rates have not moved. The supply of
negative cases has, and that is enough to change the meaning of the selected
set.

Average precision aggregates precision as the ranking is traversed, so ordinary
empirical AP evaluates performance using the positive fraction in the
evaluation data. Let $D$ denote those data: the observed outcomes and model
scores on a common set of evaluation units. When $D$ represents one intended
population, its prevalence can remain implicit. Reporting AP at the
evaluation's own prevalence answers a question about that sample. Resampling to
a target ratio discards data while still fixing a single population. Reporting
the precision--recall curve defers the choice to a reader who lacks the
sampling information needed to make it. Each leaves the claim indexed to one
population. A first task may therefore be to construct a target-indexed
quantity, $\AP_\pi(D)$, that allows for calculation of average precision at any
chosen prevalence using all of the available data. Such a quantity carries the
selection rates observed in the evaluation over to the target population
unchanged, which is a claim about that population rather than a property of the
arithmetic.

One target-prevalence value does not render a performance claim applicable to
several different deployment populations. Any single convenient prevalence
hides where performance weakens, while averaging over a distribution of target
prevalences both requires that distribution to be defensible and permits a
strong result at one prevalence to offset a weak result elsewhere. The claim
pursued here is the more demanding one: that a performance level holds
throughout a scientifically plausible prevalence range rather than in
expectation across it.

A second problem arises if the claim is intended to persist under new
evaluation units or under other sources of variation. This calls for some
notion of an evaluation regime under which a performance claim can be
repeatedly subjected to the same challenge. The extent to which a claim
survives those repeated attempts at refutation is what will be called its
support.

The population and replication problems cannot be solved independently.
Different evaluations may weaken at different prevalences, so challenging every
evaluation at one common prevalence can miss the condition that defeats each of
them. A claim required to hold throughout the population set in every
evaluation must allow the limiting prevalence to change between evaluations.

Neither problem concerns the precision of an estimate. An interval for an
expected advantage describes how far the number would move under resampling; it
does not say whether the advantage would remain in a population with a
different supply of negatives, or in the next evaluation. What is at issue is
survival rather than precision: whether the advantage remains once the
conditions under which it could fail have been named and imposed.

Call $D$ a full evaluation when it contains the outcomes and scores needed to
evaluate every declared condition on one set of units, so that no condition
requires a separate dataset. Let
\[
\Assess=(\CSet,\Regime,K),
\]
where $\CSet$ contains the conditions that can be evaluated within each full
evaluation, $\Regime$ describes how full evaluations are generated, and $K$ is
the order of the replication challenge. For AP, target prevalence will supply
the condition set; mechanisms that cannot be recovered by changing prevalence
must instead vary through $\Regime$. The order $K$ is the number of independent
opportunities to fail that the claim is required to survive. It is neither
estimated from data nor equipped with a default, because it does not describe
the data at all: it states how severe a test the claimant has chosen to face.
These elements belong to the claim rather than to the analysis. Like a
preregistered analysis plan, they lose their force if they are chosen after the
comparison has been seen.

What can then be estimated depends on the evidence available. Resampling a
single evaluation can stand in for repeated evaluation, but only under an
empirical regime closed to whatever that evaluation did not contain. The
resulting support is projected rather than observed, and no number of resamples
converts one evaluation into new empirical tests.

The motivating application we develop is model replacement: whether a candidate
model $A$ should replace an incumbent model $B$ on ranking performance. Pairing
introduces a third difficulty. Subtracting two separately challenged reports
lets each model meet its adverse condition in a different population or a
different evaluation. Forming the difference on the same units removes that
latitude at the cost of a new one, since finite-sample AP can reward whichever
model divides tied scores more finely, whether or not the division tracks the
outcomes.


\section{Average precision across target prevalences}
\label{sec:prevalence}

\subsection{Prior-standardized average precision}

Changing the observed positive fraction can alter both the population represented and the evidence used to describe the ranking. Prior standardization separates these operations by asking a counterfactual question: what precision would the observed class-conditional ranking behavior produce under a different class prior?

With $S$ denoting either fixed model score, let the observed class-conditional selection rates at threshold $\tau$ be
\[
r(\tau)=\Prb_{\mathrm{obs}}(S\geq \tau\mid Y=1),
\qquad
f(\tau)=\Prb_{\mathrm{obs}}(S\geq \tau\mid Y=0).
\]
Holding these rates fixed while varying the target prevalence is warranted only when the class-conditional score behavior transports. Because replacement is paired, the required condition for each target population is the joint invariance
\begin{equation}
\Prb_{\mathrm{obs}}\{(S_A,S_B)\in G\mid Y=y\}
=
\Prb_{\mathrm{target}}\{(S_A,S_B)\in G\mid Y=y\},
\qquad y\in\{0,1\},
\label{eq:score-invariance}
\end{equation}
for every measurable $G\subseteq\mathbb R^2$. Its marginal implications identify each model's target AP profile, while the joint condition preserves the within-class score dependence needed for the sampling law of the paired difference. This score-level form of prior-probability shift is weaker than invariance of $P(X\mid Y)$ but follows from it for fixed scoring rules \cite{saerens2002adjusting,storkey2009training,lipton2018detecting}.

Under this condition, Bayes' rule expresses precision at threshold $\tau$ in a target population with prevalence $\pi\in(0,1)$ as
\begin{equation}
p_\pi(\tau)
=
\frac{\pi r(\tau)}
{\pi r(\tau)+(1-\pi)f(\tau)}.
\label{eq:reference-precision}
\end{equation}
Its odds separate the target prior from the ranking behavior:
\begin{equation}
\frac{p_\pi(\tau)}{1-p_\pi(\tau)}
=
\frac{\pi}{1-\pi}\frac{r(\tau)}{f(\tau)}.
\label{eq:precision-odds}
\end{equation}
The functions $r(\tau)$ and $f(\tau)$ are therefore held at the values supplied by evaluation $D$ while $\pi$ indexes the declared target populations. Varying $\pi$ changes only the class prior. Aggregating $p_\pi(\tau)$ across the distinct score thresholds using the observed positive-class recall increments gives $\AP_\pi(D)$, prior-standardized AP for evaluation $D$ at prevalence $\pi$ \cite{elkan2001foundations,siblini2020calibration}. Appendix~\ref{app:ap} gives the empirical definition.

The restriction matters because prevalence can accompany other population changes. A referral population may contain more severe cases within each outcome class, while a new measurement process may change the score for otherwise similar cases. In either setting, Equation~\eqref{eq:reference-precision} reweights rates that need not describe the target. Those changes must be represented by observed evaluations from the corresponding environments rather than by prior adjustment alone \cite{bareinboim2013transportability}.

\subsection{A common effect scale}

Prior-standardized AP still has a different random-ranking value at every target prevalence. Population random ranking has $r(\tau)=f(\tau)$ and hence $\AP_\pi=\pi$. Define
\begin{equation}
\CNAP_\pi(D)
=
\frac{\AP_\pi(D)-\pi}{1-\pi}.
\label{eq:cnap}
\end{equation}
CNAP is zero for population random ranking, one for perfect ranking, and negative below chance. The normalization has the chance-corrected form $(\text{observed}-\text{chance})/(\text{ceiling}-\text{chance})$ \cite{cohen1960coefficient}. For two models,
\begin{equation}
\CNAP_{A,\pi}-\CNAP_{B,\pi}
=
\frac{\AP_{A,\pi}-\AP_{B,\pi}}{1-\pi}.
\label{eq:cnap-difference}
\end{equation}
The common chance term cancels, while division by the random-to-perfect headroom $1-\pi$ expresses the AP gap on a common opportunity scale. Thus a margin $d$ requires a raw AP gain of $d(1-\pi)$ at prevalence $\pi$; the same identity applies after tie averaging. This scaling is substantive and is appropriate when the replacement margin refers to a fraction of random-to-perfect improvement rather than an equal raw AP increment at every prevalence.

Common endpoints do not make the profile prevalence-invariant. At a threshold with true-positive and false-positive rates $r_h$ and $f_h$, the normalized contribution is
\begin{equation}
\frac{p_{\pi,h}-\pi}{1-\pi}
=
\frac{\pi(r_h-f_h)}
{\pi r_h+(1-\pi)f_h}.
\label{eq:threshold-cnap}
\end{equation}
Rare-positive populations therefore charge false positives more heavily. Two models can exchange order as $\pi$ changes even when both profiles improve with prevalence.

Let $\PSet\subset(0,1)$ be a nonempty compact set chosen from deployment plans or scientific knowledge before the paired-difference profile is examined. The CNAP profile then supplies the AP-specific condition index: $c=\pi$ and $\CSet=\PSet$. No prevalence has yet been selected or averaged away. The next problem is to determine what a claim throughout this profile must retain when the evaluation is replicated.

\section{Support across declared conditions and evaluations}
\label{sec:support}

The prevalence profile in Section~\ref{sec:prevalence} is one instance of a condition-indexed effect. More generally, let $\Delta_c(D)$ denote an effect under condition $c\in\CSet$, oriented so that larger values favor the claim. Every condition in the nonempty set $\CSet$ must be evaluable from one complete evaluation $D$. For AP under score-level prior invariance, $c$ is a target prevalence and $\CSet=\PSet$.

Not every condition can be placed in $\CSet$. When a change alters class-conditional ranking behavior, one evaluation cannot recover its effect by reweighting. Such variation must be observed through complete evaluations generated by $\Regime$. The assessment introduced above therefore separates conditions challenged within each evaluation from mechanisms challenged by generating evaluations.

\subsection{The evaluation-generating regime}

Let $\Regime$ denote the probability law for one complete evaluation $D$. A replication is one draw from that law. A newly observed draw is an \emph{empirical replication}; a bootstrap evaluation drawn from an empirical approximation $\widehat{\Regime}$ is a \emph{computational replication} intended to approximate such a draw. In either case, the replication is the complete evaluation, not the act of resampling.

Let $D_1,\ldots,D_K\iid\Regime$, with $K\geq1$. The value of $K$ belongs to the claim: it sets the order of the replication challenge and hence its lower-tail emphasis, not how many evaluations happen to be available for estimation. Independence is required between complete evaluations, although observations within an evaluation may be clustered or matched when the design and analysis preserve that structure.

A fixed-model regime draws new evaluation units while holding the fitted models, environment, and measurement process fixed. When only one complete evaluation is observed, resampling its units can approximate that regime; it does not define the regime or release mechanisms that the evaluation held fixed. A broader regime may vary environments or repeat development, but estimating it requires corresponding empirical replications or a justified generating hierarchy.

Several empirical replications do not by themselves determine how broad the regime is. They may share one measurement process or arise from a regime that releases additional mechanisms. A bootstrap can likewise redraw units from heterogeneous data without representing new heterogeneity. Scientific reach is determined by the declared regime and the adequacy of its empirical approximation, not by the number of draws alone.

\subsection{Joint survival}

Expected performance is weaker than performance required to persist. If two independent replications yield effects $0.70$ and $0.05$, their mean is $0.375$, although a claim at that level failed in the second replication. The greatest level retained by both is their minimum.

\begin{definition}[Replication-survival support]
For $D_1,\ldots,D_K\iid\Regime$, let $T_j=T(D_j)$ be an effect oriented so that larger values are better. Define
\begin{equation}
S_K(T;\Regime)
=
\E_{\Regime}\left[
\min_{1\leq j\leq K}T_j
\right].
\label{eq:support}
\end{equation}
\end{definition}

The realized minimum is the greatest level retained by one set of $K$ replications. Its expectation is an integrated joint-survival magnitude on the scale of the effect; it is neither a confidence bound nor the probability that a particular margin survives. For a prespecified margin $d$, literal joint survival is instead
\[
\psi_{K,d}
=
\Prb_{\Regime}\left(\min_{j\leq K}T_j>d\right)
=
\Prb_{\Regime}(T>d)^K.
\]
Appendix~\ref{app:literal-survival} shows that $S_K$ integrates $\psi_{K,d}$ over effect levels and gives its estimator. At $K=2$,
\begin{equation}
S_2(T;\Regime)
=
\E_{\Regime}[T]
-
\frac12\E_{\Regime}|T_1-T_2|.
\label{eq:k2}
\end{equation}
The criterion is expected performance minus half the expected disagreement between two independent replications.

The interpretation is Popperian only in a limited, negative sense. A level is retained when it survives specified opportunities for defeat, but survival does not establish truth \cite{popper1959logic,popper1963conjectures}. The regime determines the kind of failure presented, while $K$ determines how many independent opportunities are required. Neither choice is supplied by the philosophy.

\subsection{Projected and observed support}
\label{sec:projected}

One observed evaluation has not empirically faced $K$ replications. Resampling it generates computational replications from $\widehat{\Regime}$ and estimates replication-survival support under that empirical approximation to $\Regime$. The resulting quantity is \emph{projected support}. Because $\widehat{\Regime}$ is closed to score values and failure modes absent from the observed evaluation, projected support offers a narrower challenge than new empirical replications can provide.

Independent empirical replications drawn under $\Regime$ can reveal failures that the regime permits but the original evaluation did not contain. An evaluation involving a mechanism excluded from $\Regime$ has a different role: it may show that the declared assessment was too narrow rather than estimate the same criterion. When relevant empirical replications are available, they enter the estimator directly rather than serving only as an aspiration for future work.

Computational and empirical replication must consequently remain distinct. Increasing the bootstrap count $B$ reduces Monte Carlo error in the same projected quantity. Increasing $K$ changes the criterion by shifting weight toward worse evaluations. Increasing the number $M$ of observed empirical replications changes the evidence available to estimate the criterion.

\subsection{The order of the challenges}

Let
\[
Z_{\CSet}(D)
=
\inf_{c\in\CSet}\Delta_c(D).
\]
The supported advantage is
\begin{equation}
S_{K,\CSet}(\Delta;\Regime)
=
\E_{\Regime}
\left[
\min_{1\leq j\leq K}
\inf_{c\in\CSet}\Delta_c(D_j)
\right].
\label{eq:robust-support}
\end{equation}
At level $s$, the event inside the construction requires $\Delta_c(D_j)>s$ for every $c\in\CSet$ and every $j\leq K$.

Moving the condition search outside support weakens that requirement. The alternative
\[
\inf_{c\in\CSet}S_K(\Delta_c;\Regime)
\]
challenges every evaluation at the same condition before taking the infimum. In the AP application, it can therefore miss a failure at $\pi=0.01$ in one evaluation and a failure at $\pi=0.40$ in another. The intended order satisfies
\begin{equation}
S_K\left(
\inf_{c\in\CSet}\Delta_c;\Regime
\right)
\leq
\inf_{c\in\CSet}
S_K(\Delta_c;\Regime).
\label{eq:order}
\end{equation}
Equality holds when the same condition minimizes the relevant profile almost surely; Appendix~\ref{app:support-proofs} gives the proof.

The challenge becomes stronger when either index set grows.

\begin{proposition}[Challenge monotonicity]
\label{prop:monotonicity}
If $\mathcal C_1\subseteq\mathcal C_2$ and $K_1\leq K_2$, then
\[
S_{K_2,\mathcal C_2}(\Delta;\Regime)
\leq
S_{K_1,\mathcal C_1}(\Delta;\Regime).
\]
\end{proposition}

No comparable ordering holds for arbitrary changes of $\Regime$, because releasing environments or development runs changes the law of the effect rather than enlarging an index set within the same law.

\section{Constructing the paired AP effect}
\label{sec:paired}

\subsection{Marginal summaries do not compose}

A replacement decision concerns an advantage realized under the same conditions. For candidate $A$ and incumbent $B$, let
\[
A_{\PSet}(D)
=
\inf_{\pi\in\PSet}\CNAP_{A,\pi}(D),
\qquad
B_{\PSet}(D)
=
\inf_{\pi\in\PSet}\CNAP_{B,\pi}(D).
\]
Subtracting their supported marginal values allows the two models to meet different adverse prevalences and different adverse evaluations.

The error is one-sided:
\begin{equation}
S_K(A_{\PSet};\Regime)
-
S_K(B_{\PSet};\Regime)
\geq
S_K\left(
\inf_{\pi\in\PSet}
\{\CNAP_{A,\pi}-\CNAP_{B,\pi}\};
\Regime
\right).
\label{eq:noncomposition}
\end{equation}
The proof in Appendix~\ref{app:paired-proofs} uses only properties of the infimum and minimum. Separate reports therefore overstate the supported paired advantage by an amount they do not identify.

\subsection{Ordinary pairing rewards arbitrary refinement}

Pairing aligns the units and outcomes, but ordinary finite-sample AP still depends on the score vectors' tie structures. Stepwise AP is upward displaced under random ordering because it evaluates precision when positive cases enter the ranking \cite{bestgen2015random}. A fully tied score vector has one threshold and returns the population random-ranking value, while a vector of distinct scores creates more opportunities for favorable positive entries.

Four units show why this matters for replacement. Let model $A$ assign four distinct scores and model $B$ assign one common score. With two positive labels, two negative labels, and $\pi=1/2$, assign the labels uniformly over the six possibilities. Neither score vector is associated with the outcomes. Under the ordinary threshold-block convention, model $A$ has expected AP $49/72$, while model $B$ has AP $1/2$ under every assignment. Hence
\[
\E[\CNAP_A-\CNAP_B]
=
\frac{13}{36}.
\]
For two independent assignments, the supported difference remains positive at $29/216$. Pairing has favored the model that broke ties without information.

A null anchor cannot be subtracted uniformly because the displacement depends on model performance. It is large near no association but vanishes at perfect separation, where both models reach the ceiling. Appendix~\ref{app:paired-proofs} gives an exact four-unit example in which centering by the joint-permutation anchor favors the coarser model even though both models separate the classes perfectly. The defect must therefore be removed before the nonlinear population and replication challenges are applied.

\subsection{Tie-averaged average precision}
\label{sec:tie-average}

Let $\mathcal O(D)$ be the strict rankings consistent with a score vector, obtained by ordering tied units in every possible way while preserving the order of the score blocks.

\begin{definition}[Tie-averaged average precision]
\begin{equation}
\APta_\pi(D)
=
\frac{1}{|\mathcal O(D)|}
\sum_{o\in\mathcal O(D)}
\AP_\pi(D_o),
\label{eq:tie-average}
\end{equation}
where $D_o$ is evaluation $D$ resolved into strict ranking $o$. Define $\CNAPta_\pi$ through Equation~\eqref{eq:cnap}.
\end{definition}

The convention corresponds to a declared random or otherwise outcome-blind resolution of units within each score block. It is appropriate when an operational cutoff may split a block and the model supplies no within-block information. If deployment selects whole score blocks or uses a prespecified informative secondary rule, threshold-block AP or the resulting ordered ranking may be more faithful. Under the random-resolution regime, tie averaging is equivalent to adding independent continuous jitter within blocks and averaging ordinary AP; Appendix~\ref{app:tie-computation} evaluates that average without enumeration or random jitter.

Tie averaging still rewards informative refinement. A finer score improves realized AP when it orders cases correctly within a coarse block and loses AP when it orders them incorrectly. Only refinement unrelated to the outcomes receives zero expected credit.

\begin{proposition}[Neutrality to arbitrary refinement]
\label{prop:refinement}
Suppose model $A$ preserves the order of model $B$'s score blocks and, conditionally on $S_B$ and $Y$, its strict ranking is uniform over the rankings consistent with $S_B$. Then, for every $\pi$,
\[
\E\!\left[
\APta_{A,\pi}\mid S_B,Y
\right]
=
\APta_{B,\pi}.
\]
\end{proposition}

The proposition conditions on the realized outcomes, so it applies when the coarse model is informative. It removes expected credit for arbitrary refinement without requiring exchangeable evaluation units.

\subsection{The supported paired advantage}

The paired profile is
\[
\Dta_\pi(D)
=
\CNAPta_{A,\pi}(D)
-
\CNAPta_{B,\pi}(D),
\]
and the advantage retained within evaluation $D$ is
\[
\RobDta(D)
=
\inf_{\pi\in\PSet}\Dta_\pi(D).
\]
The model-replacement estimand is
\begin{equation}
\theta_{A:B}
=
S_K(\RobDta;\Regime)
=
\E_{\Regime}
\left[
\min_{1\leq j\leq K}
\inf_{\pi\in\PSet}
\{\CNAPta_{A,\pi}(D_j)-\CNAPta_{B,\pi}(D_j)\}
\right].
\label{eq:paired-estimand}
\end{equation}
The difference is formed before either minimum. A positive value is therefore a positive expected worst-of-$K$ advantage after every evaluation has been challenged throughout the declared prevalence set.

\subsection{Reference behavior and design scope}
\label{sec:paired-scope}

Refinement neutrality compares a fine ranking with a coarser ranking it refines. Two arbitrary uninformative rankings require a stronger reference law.

\begin{definition}[Exchangeable-label reference law]
\label{def:exchangeable}
Conditionally on both score vectors and the positive count $n_+$, every label vector containing $n_+$ positives has probability $\binom{n}{n_+}^{-1}$.
\end{definition}

Under this law, labels read along any fixed strict ranking are uniform over the same set of arrangements. Expected tie-averaged AP therefore depends on $n_+$, $n_-$, and $\pi$, but not on the score vector or its tie structure.

Let $\Regime_0$ denote the diagnostic reference regime. Under $\Regime_0$, each complete evaluation is generated by the exchangeable-label law, conditional on the score vectors and positive count fixed by the design. This regime diagnoses the comparison's anchor; it is not the scientific replication regime.

\begin{proposition}[Reference behavior of paired support]
\label{prop:paired-reference}
Under the exchangeable-label reference law,
\[
\E_0[\Dta_\pi]=0
\qquad\text{for every }\pi\in\PSet,
\]
and consequently
\[
S_K(\RobDta;\Regime_0)\leq0.
\]
The same bound holds after exchanging $A$ and $B$.
\end{proposition}

The inequality need not be equality because the prevalence and replication minima can move an uninformative paired effect below zero. The result prevents differing tie structures from manufacturing a positive supported advantage, although the challenge itself may be conservative.

Full exchangeability is stronger than independence. Matched, clustered, and heterogeneous-risk designs usually permit only a restricted group of label arrangements, under which score vectors aligned with the strata can have different reference values. Tie averaging still removes arbitrary ordering within score blocks, and Equation~\eqref{eq:paired-estimand} remains descriptive, but the no-positive-anchor result no longer follows. Appendix~\ref{app:paired-proofs} gives a matched-pair counterexample.

\subsection{No verdict and replacement}

The reversed orientation is
\[
\theta_{B:A}
=
S_K\left(
\inf_{\pi\in\PSet}\{-\Dta_\pi\};
\Regime
\right).
\]
Because the support operator contains minima, the reversed value is not generally the negative of the original. The two orientations nevertheless satisfy
\begin{equation}
\theta_{A:B}+\theta_{B:A}\leq0.
\label{eq:orientation}
\end{equation}
Both models therefore cannot receive a positive supported advantage over the other.

The inequality permits a no-verdict region. When both orientations are nonpositive, the declared challenges defeat superiority in either direction. The diagnostic range induced by the two orientations is
\[
\left[\theta_{A:B},-\theta_{B:A}\right].
\]
It is not an uncertainty interval. At $K=2$, its width is
\[
\E\!\left[
\max_{j\leq2}\sup_{\pi\in\PSet}\Dta_\pi(D_j)
-
\min_{j\leq2}\inf_{\pi\in\PSet}\Dta_\pi(D_j)
\right].
\]
It therefore contains both variation across target prevalences and disagreement across evaluations. When $\PSet$ is a singleton, the width reduces to the expected absolute disagreement between the paired effects of two evaluations.

Let $d\geq0$ be the prespecified smallest replication-challenged CNAP advantage that would justify replacing model $B$ by model $A$. The performance criterion is
\begin{equation}
\theta_{A:B}>d.
\label{eq:replacement}
\end{equation}
This is a risk-sensitive magnitude criterion; it does not imply a high probability that every replication clears $d$. A point estimate does not establish the inequality. The proposed decision rule requires a validated lower uncertainty bound for $\theta_{A:B}$ to exceed $d$, with pairing preserved through every estimation layer. Section~\ref{sec:estimation} constructs this bound through an outer resampling procedure.

The reverse decision requires the corresponding lower bound for $\theta_{B:A}$ to exceed its prespecified margin. If neither orientation clears its margin, the analysis returns no verdict.

\section{Evidence and estimation}
\label{sec:estimation}

\subsection{Several empirical replications}

Suppose $M$ genuinely independent empirical replications $D_1,\ldots,D_M\iid\Regime$ are available. Here $M$ counts complete evaluations; replication $D_m$ may itself contain $n_m$ evaluation units. For each replication, compute
\[
z_m
=
\inf_{c\in\CSet}\Delta_c(D_m).
\]
When $M\geq K$, the complete order-$K$ U-statistic \cite{hoeffding1948ustat}
\begin{equation}
\widehat S_{K,\CSet,M}^{\,\mathrm{obs}}
=
\binom MK^{-1}
\sum_{\substack{I\subseteq\{1,\ldots,M\}\\|I|=K}}
\min_{m\in I}z_m
\label{eq:observed-u}
\end{equation}
is unbiased for Equation~\eqref{eq:robust-support}. Sorting avoids enumerating all subsets; Appendix~\ref{app:computation} gives the resulting formula. For AP replacement, $\CSet=\PSet$ and $\Delta_\pi=\Dta_\pi$. We denote the resulting observed-evaluation estimate of $\theta_{A:B}$ by $\widehat\theta_{A:B}^{\,\mathrm{obs}}$.

Empirical replications make the inferential basis transparent, but they do not remove the need to define $\Regime$. Several test sets from one measurement process identify a narrower regime than evaluations that repeat the development procedure. When $M$ is small, the U-statistic can also be imprecise, especially as $K$ approaches $M$ and few distinct subsets remain.

\subsection{One observed evaluation}

When only one independent evaluation is available, resampling can project Equation~\eqref{eq:robust-support} only under the empirical regime $\widehat{\Regime}$ represented by that evaluation. Each bootstrap resample is one computational replication intended to approximate an independent draw of a complete evaluation from the applicable fixed-model regime. It must preserve the pairing and the joint condition profile before it records
\[
z_b
=
\inf_{c\in\CSet}\Delta_c(D^{(b)}).
\]
The $B$ computational replications approximate the evaluation-effect distribution under $\widehat{\Regime}$. Applying the complete order-$K$ U-statistic to $z_1,\ldots,z_B$ estimates the expected minimum for groups of $K$ independent draws from that distribution. For AP replacement, we denote this projected estimate by $\widehat\theta_{A:B}^{\,\mathrm{proj}}$.

For AP, let $D_{\mathrm{obs},+}$ and $D_{\mathrm{obs},-}$ contain the positive and negative evaluation units. Computational replication $b$ samples the declared numbers of units with replacement within each class, keeps both model scores attached to every sampled unit, and evaluates the complete tie-averaged CNAP difference throughout $\PSet$. Its recorded effect is
\[
z_b
=
\inf_{\pi\in\PSet}
\{\CNAPta_{A,\pi}(D^{(b)})-\CNAPta_{B,\pi}(D^{(b)})\}.
\]

Stratification conditions the projected regime on its positive and negative counts. A same-design analysis uses the observed counts, whereas a prospective analysis uses the counts planned for future evaluations. If the claim instead allows class counts to vary, their generating law belongs in $\Regime$. Clustered or matched studies must resample the coarsest unit treated as independent, which may preclude simple class-stratified sampling.

The same sampled units must be used for both models and throughout the prevalence profile. Resampling models separately would destroy pairing, while resampling separately at each prevalence would prevent one evaluation from revealing its own limiting population.

The bootstrap represents only variation available in the observed score pools. It cannot generate unseen score values or failures caused by a new environment, measurement process, or development run. Its claim must therefore be reported as projected fixed-model support rather than as survival of $K$ observed replications.

\subsection{Tail fidelity, prevalence search, and sparse outcomes}

The role of $K$ creates an estimation tension. Let $Q_T$ denote the quantile function of $T$. From Appendix~\ref{app:support-proofs},
\[
S_K(T;\Regime)
=
\int_0^1Q_T(v)K(1-v)^{K-1}\,dv,
\]
so increasing $K$ moves weight toward the lower tail of the replication-effect distribution. The stronger challenge consequently depends more heavily on whether the empirical replications, or their bootstrap approximation, represent adverse outcomes. This dependence is especially consequential when positive outcomes are sparse and the distribution of AP is coarse.

A finite set $\PSet=\{\pi_1,\ldots,\pi_G\}$ defines an exact finite challenge. A grid used to approximate an interval can miss an interior minimum and raise the reported value. A one-dimensional optimizer can instead search a compact interval, with its tolerance checked against a dense grid.

Nonunique or moving prevalence minimizers require a further distinction. Uniformly accurate profiles may still yield an accurate infimum value even when the minimizing prevalence is unstable. Inference about that value is less regular, however, because the infimum is nonsmooth at flat or multiple minima. Consistency of the projected point estimate, fidelity of its lower tail, and coverage of an outer bound must therefore be studied separately.

\subsection{Outer uncertainty}

An outer resampling layer estimates uncertainty from the observed evidence. With several empirical replications, an outer sample resamples those replications and recomputes the order-$K$ estimator. With one evaluation, each outer sample first resamples units according to the study design; an inner layer then generates computational replications and estimates projected support. The outer layer therefore repeats the entire applicable estimator. Every outer sample must preserve pairing and repeat tie averaging, the prevalence search, and the order-$K$ calculation.

Ordinary bootstrap coverage is not automatic because AP can be discrete and the minimizing prevalence can change under perturbation. Percentile, basic, studentized, and bias-corrected bounds may behave differently for the same design. The replacement rule therefore treats the lower bound as a procedure requiring design-specific simulation or theory rather than assigning nominal coverage by convention. With one observed evaluation, this bound addresses sampling uncertainty in the projected procedure; it does not account for failure modes omitted from $\widehat{\Regime}$.

\section{Interpretation, reporting, and scope}
\label{sec:discussion}

\subsection{What the number supports}

The generic estimand is an assessment-relative integrated joint-survival magnitude, not a portable property of either model. It is the expected minimum, among $K$ evaluations generated under $\Regime$, of the paired advantage after challenge throughout $\CSet$. In the AP application, $\CSet=\PSet$ and the effect is the tie-averaged CNAP difference. Omitting any of those declarations changes what the number means.

Projected and observed estimates require different language. An estimate based on several independent empirical replications may be described as replicated evidence under the declared regime. A one-evaluation bootstrap estimates projected support under empirical within-class score distributions and should not be described as having survived external replications.

The result also remains a ranking-performance claim. A deployment decision may turn on calibration or on the consequences of acting on the selected cases, neither of which Equation~\eqref{eq:replacement} evaluates.

\subsection{Minimum reporting content}

A concise report should identify the claim before presenting diagnostics. Table~\ref{tab:reporting} gives the minimum content.

\begin{table}[ht]
\centering
\caption{Minimum reporting content for a supported replacement analysis.}
\label{tab:reporting}
\begin{tabularx}{\textwidth}{@{}p{0.23\textwidth}X@{}}
\toprule
Component & Report \\
\midrule
Supported quantity &
$\widehat\theta_{A:B}$ (identified as observed or projected), the reversed orientation, the margin $d$, and the outer lower bound used for the decision \\
Assessment &
$\CSet$ and its scientific basis, with $\CSet=\PSet$ for AP; $\Regime$; what is fixed and redrawn; the independent evaluation unit; and $K$ \\
Evidence &
Whether complete evaluations are observed or bootstrap-projected; observed class counts; replication counts; and resampling hierarchy \\
Condition transport &
The invariance used to evaluate every $c\in\CSet$ from one evaluation; for AP, the score-level prior-shift assumption and its target populations \\
Comparison &
Model orientation, paired sampling, tie-averaged convention, and the basis for the exchangeable-label reference law \\
Diagnostics &
Paired-difference profile, distribution of minimizing conditions, evaluation-effect distribution, literal-margin survival $\widehat\psi_{K,d}$ (identified as observed or projected), and sensitivity to $K$ and $\CSet$ \\
Computation &
Condition set or optimizer, bootstrap count where applicable, Monte Carlo error, and interval method \\
\bottomrule
\end{tabularx}
\end{table}

A possible headline statement is:
\begin{quote}
Across the declared prevalence set $\PSet$ and under an order-$K$ replication challenge in the fixed-model regime, model $A$ had an estimated tie-averaged supported CNAP advantage of [value] over model $B$; the reversed orientation was [value]. The literal-margin survival estimate was [value]. The outer lower bound [did or did not] exceed the prespecified replacement margin $d=[value].$
\end{quote}
When the estimate comes from one test set, the opening should instead read, ``Under a bootstrap approximation to the empirical fixed-model regime, model $A$ had an estimated advantage under the prevalence and order-$K$ replication challenges.''

\subsection{Limits of the AP claim}

Prior standardization does not transport ranking behavior across changes within outcome class. If severity, measurement, or feature distributions change $P(S\mid Y)$, target-population evaluations are required.

Tie averaging does not create full label exchangeability. In matched, clustered, or heterogeneous-risk designs, the paired estimand remains computable but may inherit a design-specific anchor unless the admissible comparisons respect the design's symmetry.

The choices of $\PSet$, $\Regime$, and $K$ remain scientific commitments. Widening $\PSet$ or increasing $K$ mechanically lowers the criterion, but neither operation proves that the resulting assessment captures the failures relevant to deployment.

\subsection{Beyond prevalence: counterfactual case mix for AUROC}
\label{sec:auroc-outlook}

The support construction is not confined to prevalence or AP. Pure prior shift leaves the area under the receiver operating characteristic curve (AUROC) unchanged, but counterfactual changes in the mix of cases within each outcome class can supply an analogous condition set. Such reweightings remain limited by the case types represented in the evaluation data, so computational replications can expose only failures available within those empirical conditional distributions. Appendix~\ref{app:auroc} develops the resulting quantity, its finite-sample interpretation, and the design requirements that remain unresolved.

\section{Conclusion}

A model-replacement claim based on average precision must survive more than a favorable marginal summary. The population must be named because AP changes with prevalence, while the evaluations must be named because an expected advantage can conceal a failure to persist. These are instances of the same noncompensation requirement, but they enter the assessment differently: target prevalence can be evaluated within each replication, whereas other variation must be generated through $\Regime$.

Prior standardization evaluates each model at the target prevalences when class-conditional score behavior is invariant. Chance normalization places those values on a common population random-to-perfect scale. Tie averaging then removes expected credit for arbitrary refinement before the models are compared on the same units. The prevalence infimum and replication minimum are applied only after that paired difference has been formed:
\[
\theta_{A:B}
=
\E_{\Regime}
\left[
\min_{j\leq K}
\inf_{\pi\in\PSet}
\{\CNAPta_{A,\pi}(D_j)-\CNAPta_{B,\pi}(D_j)\}
\right].
\]

Several independent empirical replications estimate this criterion directly through a U-statistic. One evaluation can support only a bootstrap projection under an empirical fixed-model regime. That projection can guide a replacement analysis, but it cannot substitute computational resampling for new empirical opportunities to reveal failure.

The condition-indexed support operator is not confined to prevalence. Counterfactual AUROC case mixes provide one possible extension, but their admissible weights and empirical resolution require a separate method before they can carry the same replacement interpretation.

\appendix

\section{Empirical prior-standardized average precision}
\label{app:ap}

\subsection{Definition and weighted implementation}

Let
\[
D=\{(s_i,y_i)\}_{i=1}^{n},
\qquad
y_i\in\{0,1\},
\]
with $n_+=\sum_i y_i>0$ and $n_-=n-n_+>0$. Let $c_1>\cdots>c_H$ be the distinct observed score values. At threshold $h$, define
\[
\operatorname{TP}_h
=
\sum_i\mathbf 1\{y_i=1,s_i\geq c_h\},
\qquad
\operatorname{FP}_h
=
\sum_i\mathbf 1\{y_i=0,s_i\geq c_h\},
\]
and
\[
r_h=\frac{\operatorname{TP}_h}{n_+},
\qquad
f_h=\frac{\operatorname{FP}_h}{n_-}.
\]
With $r_0=0$, noninterpolated prior-standardized AP under the threshold-block convention is
\begin{equation}
\AP_\pi(D)
=
\sum_{h=1}^{H}
(r_h-r_{h-1})
\frac{\pi r_h}
{\pi r_h+(1-\pi)f_h}.
\label{eq:empirical-ap}
\end{equation}
All units sharing a score enter one threshold block. A block containing only negative cases has zero recall increment, although its false positives reduce the precision of later blocks.

The same quantity can be computed by weighting each positive unit by $w_+=\pi/n_+$ and each negative unit by $w_-=(1-\pi)/n_-$. Weighted precision at threshold $h$ is Equation~\eqref{eq:reference-precision}, while weighted recall remains $r_h$.

\subsection{Behavior across prevalence}

Combining Equations~\eqref{eq:cnap} and \eqref{eq:empirical-ap} gives
\begin{equation}
\CNAP_\pi(D)
=
\sum_{h=1}^{H}
(r_h-r_{h-1})
\frac{\pi(r_h-f_h)}
{\pi r_h+(1-\pi)f_h}.
\label{eq:cnap-profile}
\end{equation}
The derivative of a threshold contribution is
\begin{equation}
\frac{\partial}{\partial\pi}
\left\{
\frac{\pi(r_h-f_h)}
{\pi r_h+(1-\pi)f_h}
\right\}
=
\frac{f_h(r_h-f_h)}
{\{\pi r_h+(1-\pi)f_h\}^2}.
\label{eq:cnap-derivative}
\end{equation}
If $r_h\geq f_h$ at every threshold with positive recall increment, the profile is nondecreasing and its minimum over $[\pi_L,\pi_U]$ occurs at $\pi_L$. A paired difference can remain nonmonotone even when both marginal profiles are nondecreasing.

\section{Properties of replication-survival support}
\label{app:support-proofs}

\subsection{Survival and quantile representations}

For a random variable $T\in[L,U]$,
\[
T=L+\int_L^U\mathbf 1\{T>s\}\,ds.
\]
For $K$ independent copies, expectation of their minimum gives
\begin{equation}
S_K(T;\Regime)
=
L+\int_L^U\Prb_{\Regime}(T>s)^K\,ds.
\label{eq:survival-representation}
\end{equation}
At each level $s$, the powered survival probability is the probability that all $K$ evaluations retain the level.

\subsection{Literal survival of a practical margin}
\label{app:literal-survival}

For a prespecified practical margin $d$, define
\begin{equation}
\psi_{K,d}
=
\Prb_{\Regime}\left(\min_{j\leq K}T_j>d\right)
=
\Prb_{\Regime}(T>d)^K.
\label{eq:literal-survival}
\end{equation}
Unlike $S_K$, this quantity is a probability at one effect level. Equation~\eqref{eq:survival-representation} can be written
\[
S_K(T;\Regime)=L+\int_L^U\psi_{K,s}\,ds,
\]
so replication-survival support is an integrated joint-survival magnitude, whereas $\psi_{K,d}$ is literal survival of the decision margin.

For independent empirical-replication effects $z_1,\ldots,z_M$, let $R_d=\sum_{m=1}^M\mathbf 1\{z_m>d\}$. When $M\geq K$, the complete order-$K$ estimator is
\begin{equation}
\widehat\psi_{K,d}
=
\binom MK^{-1}
\sum_{\substack{I\subseteq\{1,\ldots,M\}\\|I|=K}}
\prod_{m\in I}\mathbf 1\{z_m>d\}
=
\frac{\binom{R_d}{K}}{\binom MK},
\label{eq:literal-survival-estimator}
\end{equation}
with $\binom{R_d}{K}=0$ when $R_d<K$. It is unbiased for $\psi_{K,d}$; applied to computational replications, it estimates survival only under the projected empirical regime. For small $M$, the estimator is coarse and equals zero whenever fewer than $K$ empirical replications clear $d$, reflecting that no observed order-$K$ subset survived. A probability-based replacement rule would additionally require a prespecified probability threshold and a validated lower bound. Because AP is discrete, the choice between $>$ and $\geq$ and sensitivity to mass at $d$ must be declared.

\subsection{Quantile representation}

If $Q_T$ is the quantile function of $T$, the smallest of $K$ independent uniform variables has density $K(1-v)^{K-1}$. Hence
\begin{equation}
S_K(T;\Regime)
=
\int_0^1Q_T(v)K(1-v)^{K-1}\,dv.
\label{eq:quantile-weight}
\end{equation}
The weight moves toward the lower tail as $K$ increases.

\subsection{Proof of the operator-order inequality}

Let $X_c^{(j)}=\Delta_c(D_j)$ and $Y_c=\min_{j\leq K}X_c^{(j)}$. For every realized array,
\[
\min_j\inf_{c\in\CSet}X_c^{(j)}
=
\inf_{c\in\CSet}\min_jX_c^{(j)}
=
\inf_{c\in\CSet}Y_c.
\]
Since $\inf_c Y_c\leq Y_{c'}$ for every $c'\in\CSet$,
\[
\E\left[\inf_{c\in\CSet}Y_c\right]
\leq
\inf_{c\in\CSet}\E[Y_c],
\]
which proves Equation~\eqref{eq:order}.

\subsection{Proof of challenge monotonicity}

Take $\mathcal C_1\subseteq\mathcal C_2$ and $K_1\leq K_2$. Draw $K_2$ evaluations once and use the first $K_1$ for the smaller assessment. For every realized array,
\[
\min_{j\leq K_2}\inf_{c\in\mathcal C_2}\Delta_c(D_j)
\leq
\min_{j\leq K_1}\inf_{c\in\mathcal C_1}\Delta_c(D_j).
\]
Expectation preserves the inequality and proves Proposition~\ref{prop:monotonicity}.

\section{Paired comparison proofs and counterexamples}
\label{app:paired-proofs}

\subsection{Why marginal subtraction overstates paired advantage}

Let $U_\pi$ and $V_\pi$ be two profiles, with
\[
U_{\PSet}=\inf_{\pi\in\PSet}U_\pi,
\qquad
V_{\PSet}=\inf_{\pi\in\PSet}V_\pi,
\qquad
T_{\PSet}=\inf_{\pi\in\PSet}(U_\pi-V_\pi).
\]
Choosing a prevalence minimizing $U_\pi$ gives
\[
T_{\PSet}
\leq
U_{\PSet}-V_\pi
\leq
U_{\PSet}-V_{\PSet}.
\]
Monotonicity of $S_K$ therefore gives
\[
S_K(U_{\PSet}-V_{\PSet})
\geq
S_K(T_{\PSet}).
\]

For effects evaluated on the same replications, the minimum is superadditive:
\[
\min_j(X_j+Y_j)
\geq
\min_jX_j+\min_jY_j.
\]
Taking $X=U_{\PSet}-V_{\PSet}$ and $Y=V_{\PSet}$ yields
\[
S_K(U_{\PSet})-S_K(V_{\PSet})
\geq
S_K(U_{\PSet}-V_{\PSet}),
\]
which completes Equation~\eqref{eq:noncomposition}.

\subsection{Proof of refinement neutrality}

Condition on $S_B$ and $Y$. Proposition~\ref{prop:refinement} makes model $A$'s strict ranking $O_A$ uniform on $\mathcal O(D_B)$. Tie-averaged AP for $A$ is the conditional expectation of ordinary AP over whatever ties $A$ retains. By the tower property,
\[
\E[\APta_{A,\pi}\mid S_B,Y]
=
\E[\AP_\pi(D_{O_A})\mid S_B,Y].
\]
Uniformity of $O_A$ turns the right side into the average of $\AP_\pi$ over $\mathcal O(D_B)$, which is $\APta_{B,\pi}$.

\subsection{Rank invariance under exchangeable labels}

Fix $n_+$, $n_-$, and a strict ranking $o$. Under Definition~\ref{def:exchangeable}, the label sequence read along $o$ is uniform over the $\binom n{n_+}$ arrangements containing $n_+$ positives. Reading along another strict ranking is a bijection of the same arrangements, so the sequence has the same law.

For a strict ranking, $\pi\mapsto\CNAP_\pi(D_o)$ is a deterministic function of that sequence. Its expectation is therefore the same for every strict ranking. Tie-averaged AP is a convex combination of strict-ranking values, so any two score vectors on the same units satisfy
\[
\E_0[\CNAPta_{A,\pi}]
=
\E_0[\CNAPta_{B,\pi}].
\]
Thus $\E_0[\Dta_\pi]=0$.

Finally,
\[
\E_0[\RobDta]
\leq
\inf_{\pi\in\PSet}\E_0[\Dta_\pi]
=0,
\]
and $S_K(T)\leq\E[T]$. Their composition proves Proposition~\ref{prop:paired-reference}; replacing $\Dta$ by $-\Dta$ proves the reversed orientation.

\subsection{Why restricted label symmetry is insufficient}

Take four units in two matched pairs, each containing exactly one positive. The design permits labels to swap within pairs but not between pairs. Let model $A$ give all units the same score, while model $B$ places the first pair above the second and ties the units within each pair. At $\pi=1/2$, exact tie averaging over the four permitted assignments gives
\[
\CNAPta_A-\CNAPta_B
=
\frac1{36}
\]
in one orientation for every assignment. The effect has no replication variation, so its supported value remains $1/36$ for every $K$.

The obstruction is the symmetry group of the design. Full exchangeability places every strict ranking in one orbit of the label-permutation group. Restricted designs may divide rankings into several orbits with different anchors, and averaging within score ties cannot move a ranking between them.

\subsection{Why paired null centering fails}

Take two positive and two negative units at $\pi=1/2$, with
\[
s_A=(4,3,2,1),
\qquad
s_B=(2,2,1,1),
\]
and let the two highest-scored units be positive. Both models separate the classes perfectly, so every stratified bootstrap evaluation has CNAP one for each model and paired difference zero.

Under joint label permutation, exact enumeration of the nested $K=2$ procedure gives a design anchor of $35/384$ for $A-B$ and $-295/1152$ for $B-A$. Subtracting these anchors would report a positive corrected value of about $0.256$ for the second orientation, favoring the coarser model although both predictions are perfect. The displacement is not an additive constant because perfect separation leaves no room for it.

\subsection{Proof of the orientation bound}

For evaluation $j$, write
\[
L_j=\inf_{\pi\in\PSet}\Dta_\pi(D_j),
\qquad
U_j=\sup_{\pi\in\PSet}\Dta_\pi(D_j).
\]
The two orientations satisfy
\[
\theta_{A:B}
=
\E\left[\min_{j\leq K}L_j\right],
\qquad
\theta_{B:A}
=
-\E\left[\max_{j\leq K}U_j\right].
\]
Since $L_j\leq U_j$ for every evaluation,
\[
\theta_{A:B}+\theta_{B:A}
=
\E\left[
\min_{j\leq K}L_j-\max_{j\leq K}U_j
\right]
\leq0,
\]
which proves Equation~\eqref{eq:orientation}. The negative of this sum is the width of the no-verdict region. When $\PSet$ is a singleton and $K=2$, $L_j=U_j=T_j$, so the width reduces to $\E|T_1-T_2|$.

\section{Closed-form computation of tie-averaged AP}
\label{app:tie-computation}

Index score blocks $j=1,\ldots,J$ from highest to lowest. Let block $j$ contain $a_j$ positive and $b_j$ negative units, with $m_j=a_j+b_j$, and define
\[
A_j=\sum_{\ell<j}a_\ell,
\qquad
B_j=\sum_{\ell<j}b_\ell.
\]
Then
\begin{equation}
\APta_\pi(D)
=
\sum_{j=1}^{J}
\frac{a_j}{n_+}\frac1{m_j}
\sum_{k=1}^{m_j}
\sum_{t=\max(0,k-1-b_j)}^{\min(k-1,a_j-1)}
\frac{\binom{a_j-1}{t}\binom{b_j}{k-1-t}}
{\binom{m_j-1}{k-1}}
\frac{\pi r_{jkt}}
{\pi r_{jkt}+(1-\pi)f_{jkt}},
\label{eq:tie-closed}
\end{equation}
where
\[
r_{jkt}=\frac{A_j+t+1}{n_+},
\qquad
f_{jkt}=\frac{B_j+k-1-t}{n_-}.
\]

To obtain the formula, fix one positive unit in block $j$. Under a uniformly random ordering of the block, its position $k$ is uniform on $\{1,\ldots,m_j\}$. Given $k$, the number $t$ of other positive units in the $k-1$ preceding positions is hypergeometric with $a_j-1$ available positives and $b_j$ negatives. The final fraction in Equation~\eqref{eq:tie-closed} is the prior-standardized precision when that positive enters the ranking. Averaging over $t$ and $k$, multiplying by the block's $a_j$ positive units, and summing over blocks gives the result.

The calculation costs $O(\sum_jm_j^2)$ operations. Its combinatorial weights do not depend on $\pi$ and can be reused throughout the prevalence search. When all scores are distinct, Equation~\eqref{eq:tie-closed} reduces to ordinary stepwise AP.

\section{Computation and outer uncertainty}
\label{app:computation}

\subsection{General support order}

Sort $M$ observed empirical-replication effects:
\[
z_{(1)}\leq\cdots\leq z_{(M)}.
\]
The value $z_{(i)}$ is the minimum in every order-$K$ subset containing it and $K-1$ of the $M-i$ larger values. Hence Equation~\eqref{eq:observed-u} equals
\begin{equation}
\widehat S_{K,\CSet,M}^{\,\mathrm{obs}}
=
\binom MK^{-1}
\sum_{i=1}^{M-K+1}
\binom{M-i}{K-1}z_{(i)}.
\label{eq:sorted-observed}
\end{equation}

For $B$ bootstrap effects, the projected estimator is
\begin{equation}
\widehat S_{K,\CSet,B}^{\,\mathrm{proj}}
=
\binom BK^{-1}
\sum_{i=1}^{B-K+1}
\binom{B-i}{K-1}z_{(i)}.
\label{eq:sorted-bootstrap}
\end{equation}
At $K=2$, this reduces to
\[
\widehat S_{2,\CSet,B}^{\,\mathrm{proj}}
=
\frac{2}{B(B-1)}
\sum_{i=1}^{B-1}(B-i)z_{(i)}.
\]

\subsection{Monte Carlo error}

For $K=2$, define the leave-one-replication projection
\[
\widehat h_{1,-i}(z_i)
=
\frac1{B-1}\sum_{j\ne i}\min(z_i,z_j).
\]
A first-order Monte Carlo standard error for the projected estimator is
\begin{equation}
\widehat{\operatorname{se}}_{\mathrm{MC}}
=
\frac2{\sqrt B}
\operatorname{sd}_{i}
\left\{
\widehat h_{1,-i}(z_i)
\right\}.
\label{eq:mcse}
\end{equation}
Prefix sums over the sorted array compute all projections in linear time after sorting.

\subsection{Prevalence evaluation and outer bounds}

For a finite prevalence set with $G$ values, score-block counts are computed once in each evaluation and reused across prevalences. For a continuous interval, Equation~\eqref{eq:cnap-profile} is a smooth one-dimensional objective away from thresholds with no recall contribution. Endpoint values and stationary points may be checked directly, while a numerical optimizer should be validated against a dense grid.

With several empirical replications, an outer sample resamples the complete evaluations. With one observed evaluation, an outer sample resamples the independent units or clusters according to the study design, and the inner bootstrap then generates computational replications from that outer empirical distribution. Pairing must be retained in both layers, and every outer sample must repeat the prevalence search. Because the target combines a lower-tail functional with a potentially nonsmooth infimum, coverage should be established under designs resembling the intended application.

\section{Counterfactual case-mix support for AUROC}
\label{app:auroc}

\subsection{Prior shift supplies no AUROC challenge}

The condition set used throughout is a set of target prevalences. It is
available because AP depends on the class prior. AUROC does not. It is the
probability that a randomly drawn positive outranks a randomly drawn negative,
with half credit for a tie, and that probability is a functional of the two
class-conditional score distributions alone. Changing the prior changes how
often each class is drawn while leaving both conditionals intact, so an AUROC
advantage that survives every $\pi\in\PSet$ has survived nothing. Whether the
support operator is a general device or an artifact of AP's prevalence
dependence therefore turns on whether AUROC admits a counterfactual of its own.

The invariance itself shows where a counterfactual is available. Prior shift leaves the composition of each class untouched, which is precisely the dimension free to differ between the evaluation and the target population even when the prior does not.
Let $X$ be a prespecified case type, and let $q_1(a)$ and $q_0(b)$ give the
target proportions of positive cases of type $a$ and negative cases of type
$b$. Writing
\[
U_{ab}
=
\Prb(S_a^+>S_b^-)
+\tfrac12\Prb(S_a^+=S_b^-)
\]
for the type-conditional comparison probability, the target case-mix AUROC is
\begin{equation}
\operatorname{AUC}_{q_1,q_0}
=
\sum_{a,b}q_1(a)\,q_0(b)\,U_{ab}.
\label{eq:auroc-mix}
\end{equation}
Prevalence reweights between the outcome classes; $(q_1,q_0)$ reweights within
them. Because $U_{ab}$ already carries half credit, exchangeable labels give
Equation~\eqref{eq:auroc-mix} mean $1/2$ under any tie structure, so the
positive mean anchor that forced tie averaging in Section~\ref{sec:tie-average}
does not recur here.

\subsection{Chosen weights and estimated components}

The weights $q_1$ and $q_0$ are
declared: they express the population the claim is about, and any nonnegative
normalized choice is admissible. The components $U_{ab}$ are not declared but
estimated, since each is a property of the score distributions that the two
named case types actually induce. A counterfactual mix therefore rescales
ingredients; it cannot supply one.

This is why knowing every observed outcome does not by itself determine a
counterfactual value. Suppose $X$ distinguishes types $A$ and $B$, and the
counterfactual asks how the models would compare if positives were
predominantly of type $B$ and negatives predominantly of type $A$. The answer
is governed by $U_{BA}$, and the labels of the cases in hand say nothing about
it unless some of those cases are type-$B$ positives and type-$A$ negatives.
What is missing when the mix cannot be evaluated is not ground truth but the
conditional score behavior of a component the evaluation never sampled.

The requirement on the data is accordingly one of coverage rather than
agreement. The evaluation need not exhibit the target proportions, exactly as a
dataset need not have prevalence $\pi$ to estimate $\AP_\pi$; it must satisfy
\begin{equation}
q_y(x)>0
\quad\Longrightarrow\quad
\text{the evaluation contains cases with }Y=y,\ X=x,
\label{eq:auroc-coverage}
\end{equation}
together with the assumption that $P(S\mid Y,X)$ transports. Coverage fails in
two distinguishable ways. If a required component is empty, the corresponding
$U_{ab}$ is undefined and the mix can be evaluated only by extrapolation. If it
holds a single case, the mix is computable, but its value rests on one score
and carries almost no information about how further cases of that type would
behave. The distinction matters because the support operator treats every
element of $\CSet$ as an opportunity for defeat, and an element whose value
cannot be reliably determined does not supply one: the challenge has been named
but not performed.

\subsection{Reweighting observed cases alone}

One construction avoids the coverage problem outright. Take the observed cases
as the universe and admit arbitrary reweightings of them, so that positives
receive weights $w_+$ and negatives weights $w_-$, each summing to one, and
\begin{equation}
\operatorname{AUC}_m(w_+,w_-)
=
\sum_{i:\,Y_i=1}\ \sum_{k:\,Y_k=0}
w_{+,i}\,w_{-,k}\,\phi(s_i^m,s_k^m),
\label{eq:auroc-empirical-mix}
\end{equation}
where $\phi$ is one for a correctly ordered pair, one half for a tie, and zero
otherwise. Every admissible mix is then exactly calculable, no transport
assumption is required, and Equation~\eqref{eq:auroc-coverage} becomes vacuous.
The construction also preserves the property one would most want preserved:
because a perfectly separating model makes every $\phi$ equal to one,
Equation~\eqref{eq:auroc-empirical-mix} equals one for every choice of weights,
so no reweighting of the evaluation can penalize a model that ranks the classes
correctly.

The difficulty is that unrestricted weights allow a mix to be supported by
arbitrarily little of the evaluation. Placing all positive weight on one
positive case and all negative weight on one negative case reduces
Equation~\eqref{eq:auroc-empirical-mix} to a single $\phi$, so the infimum over
mixes is the worst pair in the dataset and any model with one inversion is
assigned zero. The paired criterion of Equation~\eqref{eq:robust-support}
degenerates further. If candidate $A$ separates perfectly then
$\phi_A-\phi_B=1-\phi_B\geq0$ for every pair, yet the infimum is attained by
concentrating on a pair that the incumbent $B$ also orders correctly, which
gives a supported advantage of exactly zero. A perfect candidate can show a
strictly positive advantage only against an incumbent that misorders every
pair. The unrestricted construction does not make replacement demanding; it
makes it unavailable.

\subsection{Bounded concentration}

What defeats the criterion is concentration rather than finite-sample
reweighting itself, so the repair is to bound it. For $\Gamma\geq1$, admit the
mixes satisfying
\begin{equation}
0\leq w_{+,i}\leq\frac{\Gamma}{n_+},
\qquad
0\leq w_{-,k}\leq\frac{\Gamma}{n_-},
\qquad
\sum_i w_{+,i}=\sum_k w_{-,k}=1,
\label{eq:auroc-gamma}
\end{equation}
and write $\CSet_\Gamma$ for the resulting condition set. At $\Gamma=1$ only
the observed mix is admissible; as $\Gamma$ approaches $n_+$ and $n_-$ the set
grows to the degenerate worst-pair criterion. Equation~\eqref{eq:auroc-gamma}
also has a direct sampling interpretation, since it forces each class's
effective sample size $(\sum_i w_i^2)^{-1}$ to be at least $n_+/\Gamma$ or
$n_-/\Gamma$; the parameter is the factor by which a counterfactual population
is permitted to shrink the evaluation that speaks to it. Because
$\Gamma_1\leq\Gamma_2$ implies $\CSet_{\Gamma_1}\subseteq\CSet_{\Gamma_2}$,
Proposition~\ref{prop:monotonicity} applies without modification and the
supported advantage is nonincreasing in $\Gamma$.

Under Equation~\eqref{eq:auroc-gamma} a perfectly separating model still
attains one throughout $\CSet_\Gamma$, and a perfect candidate now has a
positive supported advantage over an imperfect incumbent unless the incumbent
orders correctly every pair within some admissible subpopulation. Replacement
is thereby supported when the incumbent's weakness persists across all declared
mixes, and refused when the weakness can be evaded by a population that retains
a stated fraction of each class. Choosing $\Gamma$ fixes the resolution of the
claim rather than its verdict, which is why it must be declared with $\CSet$
and $\Regime$ before the values are seen.

\subsection{What reweighting cannot reach}

Bounded concentration resolves how much of the evaluation a counterfactual
population must rest on. It does not extend the claim beyond the case types the
evaluation contains. For any $\Gamma$, the admissible mixes are recombinations
of observed cases, and a type absent from the evaluation receives weight zero in
all of them; reaching it requires new data or an explicit extrapolation model.

The boundary is the one already met in the AP construction. Prior
standardization reweights the observed class-conditional distributions and the
bootstrap resamples from them, so both operations explore recombinations of a
fixed empirical world while holding $P(S\mid Y)$ invariant. Case-mix
reweighting does the same for $P(S\mid Y,X)$. Neither can reveal an unseen case
type, score behavior outside the observed support, or a failure of the
invariance each assumes. Applying the support operator to $K$ computational
replications accordingly asks what would survive $K$ hypothetical evaluations
under an empirical approximation, whereas the same operator applied to
independent empirical evaluations measures survival against challenges able to
contradict what the projection holds fixed.

Bounded concentration therefore completes the definition without completing the
criterion. The remaining obstacle is that $\Gamma$ has no referent outside the
analysis. A prevalence set is nameable because deployment prevalence is a fact
about the setting, so $\PSet$ is fixed by something the analyst does not
control; nothing about a deployment states how far a case mix may concentrate,
and $\Gamma$ is left to be chosen by whoever can also observe its effect on the
result. The choice is consequential in two directions at once, since a larger
bound admits more mixes and so strengthens the challenge, while leaving each
admissible mix supported by fewer effective observations and so estimated less
precisely. Supplying an external standard — a characterization of the
populations a deployment can present, or a bound tying $\Gamma$ to a stated
estimation precision — is what would make Equation~\eqref{eq:robust-response}
an AUROC replacement rule rather than a family of them indexed by an
undetermined parameter.


\begin{thebibliography}{99}

\bibitem{bareinboim2013transportability}
Bareinboim, E., \& Pearl, J. (2013).
A general algorithm for deciding transportability of experimental results.
\emph{Journal of Causal Inference}, 1(1), 107--134.

\bibitem{bestgen2015random}
Bestgen, Y. (2015).
Exact expected average precision of the random baseline for system evaluation.
\emph{The Prague Bulletin of Mathematical Linguistics}, 103, 131--138.

\bibitem{cohen1960coefficient}
Cohen, J. (1960).
A coefficient of agreement for nominal scales.
\emph{Educational and Psychological Measurement}, 20, 37--46.

\bibitem{elkan2001foundations}
Elkan, C. (2001).
The foundations of cost-sensitive learning.
In \emph{Proceedings of the Seventeenth International Joint Conference on Artificial Intelligence}, 973--978.

\bibitem{hoeffding1948ustat}
Hoeffding, W. (1948).
A class of statistics with asymptotically normal distribution.
\emph{Annals of Mathematical Statistics}, 19, 293--325.

\bibitem{lipton2018detecting}
Lipton, Z. C., Wang, Y.-X., \& Smola, A. (2018).
Detecting and correcting for label shift with black box predictors.
In \emph{Proceedings of the 35th International Conference on Machine Learning}, PMLR 80, 3122--3130.

\bibitem{popper1959logic}
Popper, K. R. (1959).
\emph{The Logic of Scientific Discovery}.
Hutchinson.

\bibitem{popper1963conjectures}
Popper, K. R. (1963).
\emph{Conjectures and Refutations: The Growth of Scientific Knowledge}.
Routledge.

\bibitem{saerens2002adjusting}
Saerens, M., Latinne, M., \& Decaestecker, C. (2002).
Adjusting the outputs of a classifier to new a priori probabilities.
\emph{Neural Computation}, 14(1), 21--41.

\bibitem{siblini2020calibration}
Siblini, W., Fr\'ery, J., He-Guelton, L., Obl\'e, F., \& Wang, Y. (2020).
Master your metrics with calibration.
In \emph{Advances in Intelligent Data Analysis XVIII}, 457--469.

\bibitem{storkey2009training}
Storkey, A. (2009).
When training and test sets are different: characterizing learning transfer.
In \emph{Dataset Shift in Machine Learning}, 3--28. MIT Press.

\end{thebibliography}

\end{document}
